% M12: Optionale Zeilenumbrueche an Paketnamensbestandteilen.
% M08, 11.09.2026: Kapitelauftrag und Zuordnung konzentriert.
\chapter{Technische Realisierung mit dem Microsoft Agent Framework}
\label{chap:maf-realisierung}

Kapitel~\ref{chap:loesungsarchitektur} hat die Architektur und ihre
Verantwortungsverteilung erklärt. Dieses Kapitel zeigt, wie die
tragenden Mechanismen mit MAF und C\#/.NET realisiert wurden:
Welche Frameworkmittel tragen die Ausführung, welche fachlichen
Regeln ergänzt das Projekt, und wo liegen technische Grenzen?
Die Abschnitte vertiefen diese Fragen entlang der Zuordnung in
Tabelle~\ref{t:maf-zuordnung}.

\section{Technischer Kontext und MAF-Zuordnung}
\label{sec:technischer-kontext-maf-zuordnung}

Die Darstellung bezieht sich auf den festzuhaltenden
Implementierungsstand des Artefakts
(Stand: \emph{Tag/Commit folgen mit dem technischen Einfrieren des
Artefakts}). Die Belege in Kapitel~\ref{chap:evaluation} werden
dagegen jeweils dem tatsächlich ausgeführten Systemstand
zugeordnet; historische Läufe gelten nicht als Ausführungen eines
späteren Stands. Die MAF-Kernnutzung basiert auf
\texttt{Microsoft.\allowbreak Agents.\allowbreak AI}, \texttt{Microsoft.\allowbreak Agents.\allowbreak AI.\allowbreak Workflows}
und \texttt{Microsoft.\allowbreak Agents.\allowbreak AI.\allowbreak OpenAI}, jeweils in Version~1.15.0.
Weitere Abhängigkeiten und Konfigurationsangaben stehen bei den
betreffenden Mechanismen und in
Abschnitt~\ref{sec:realisierung-reproduzierbarkeit-grenzen}.

Die allgemeine Grenze zwischen Frameworkangebot und Projektaufgabe
aus Abschnitt~\ref{sec:maf-leistungsumfang-grenzen} wird hier am
Artefakt konkretisiert. Tabelle~\ref{t:maf-zuordnung} verbindet jede
Architekturverantwortung mit ihrer technischen Grundlage und dem
Ort ihrer Vertiefung. MAF trägt etwa Nachrichtenübergaben,
Werkzeugaufrufe und Workflowzustand; das Evidenzmodell, die
fachliche Geltung von Änderungen und die Regeln der Außenabbildung
legt das Projekt fest. Die konkrete Zuordnung richtet sich nach
dem verwendeten Code, nicht allein nach dem verfügbaren
Frameworkangebot. Dokumentierte Auswahlentscheidungen und
relevante, nicht eingesetzte Alternativen werden am jeweiligen
Mechanismus erläutert.

\input{tables/chap6/maf-zuordnung}
